% Vietnamese translation for OI Public Library
% Source-File: IOI中国国家候选队论文/2005/龙凡/龙凡.ppt
\documentclass[11pt]{article}
\usepackage{oipl-vn}

\title{Bản trình chiếu: Ứng dụng của thứ tự}
\author{Long Fan}
\date{2005}

\begin{document}
\maketitle

\origtitle{Slide companion for applications of order}
\sourcepath{IOI China candidate team papers / 2005 / Long Fan / Long Fan.ppt}

\section{Phạm vi bản dịch}

Bản trình chiếu đi cùng bài viết về quan hệ thứ tự. Phần chữ đọc được chứa
các công thức \(t(i)\), \(t(j)+1\), các độ phức tạp \(O(N\log N)\), \(O(N)\),
\(O(N^2)\), \(O(N\log(N+M))\), và các nhãn \texttt{Goto}, \texttt{NewGoto}.

\section{Tư tưởng}

Nhiều bài toán có một thứ tự ẩn: thứ tự thời gian, thứ tự topo, thứ tự giá
trị, hoặc thứ tự bao hàm. Khi tìm được thứ tự này, ta có thể thay tìm kiếm
toàn cục bằng quét tuyến tính, quy hoạch động theo thứ tự, hoặc cấu trúc dữ
liệu duy trì tiền tố.

\section{Ví dụ chuyển trạng thái}

Các nhãn \texttt{Goto} và \texttt{NewGoto} gợi ý một bài toán tự động hoặc
chuyển trạng thái. Nếu trạng thái có thể sắp theo thứ tự, mỗi chuyển tiếp mới
có thể được tính từ chuyển tiếp trước, làm giảm độ phức tạp từ \(O(N^2)\)
xuống \(O(N\log(N+M))\) hoặc \(O(N+M)\).

\section{Kết luận}

Thứ tự là một quan hệ mạnh nhưng thường bị ẩn trong đề. Bài trình bày nhấn
mạnh việc tìm ra thứ tự đúng trước khi chọn kỹ thuật cài đặt.

\end{document}
